\chapter{Tag 4 — SECDED, Bündelfehler, Interleaving, CRC}

Stub-Kapitel für M1. Inhalte folgen in M4.

Dieses Kapitel testet \emph{Szenario 2} aus dem Briefing: eine Aufgabe,
die über einen Seitenumbruch geht. Der farbige Strich am Rand muss
sich auf der nächsten Seite fortsetzen, die Marginalie soll aber nur
am Anfang stehen.

\begin{bleistiftuebung}{1}{Eine bewusst lange Aufgabe für den Seitenumbruch-Test}
Die folgende Aufgabe ist absichtlich ausführlich gehalten, damit der
Aufgaben-Block über einen Seitenumbruch hinausgeht.

Stell dir eine zerkratzte CD vor. Ein einzelner Kratzer von wenigen
Millimetern zerstört nicht ein einzelnes Bit, sondern eine ganze
Sequenz benachbarter Bits — ein \textbf{Bündelfehler} (engl.\ \emph{burst
error}). Hamming-Codes, wie wir sie an Tag 3 gesehen haben, sind genau
für diese Situation \emph{nicht} gemacht: sie korrigieren einen
einzelnen Bitfehler, lügen aber bei zwei oder mehr.

Es gibt zwei prinzipielle Antworten: stärkere Codes verwenden (das
führt uns zu Reed-Solomon ab Tag 7) oder den Bündelfehler in mehrere
Einzelfehler zerlegen, indem man die Daten vor dem Übertragen so
\emph{verteilt}, dass benachbarte Sendebits in unterschiedlichen
Codewörtern landen. Dieses Verteilen heißt \textbf{Interleaving}.

Aufgabe: nimm einen \(3 \times 4\)-Block aus zwölf Bits, schreibe ihn
zuerst zeilenweise auf, dann spaltenweise. Markiere einen Bündelfehler
über vier benachbarte Bits in der spaltenweisen Reihenfolge. Wo landen
die fehlerhaften Bits in der zeilenweisen Anordnung? Wenn jede Zeile
ein Codewort ist, das einen Einzelfehler korrigieren kann — wie viele
der vier Fehler werden korrigiert?

Diese Übung ist absichtlich textlastig. Sie zeigt, dass das Layout den
Strich auf der nächsten Seite konsequent fortsetzt und dass die
Marginalie am Aufgabenanfang sitzen bleibt. Wenn die Aufgabe nach dem
Seitenumbruch bündig weiterläuft, ist das Layout in diesem Punkt
einsatzbereit.
\end{bleistiftuebung}
